% Generated draft: requires mathematical and source review.
\documentclass[11pt]{article}
\usepackage[margin=25mm]{geometry}
\usepackage{fontspec}
\setmainfont{Latin Modern Roman}
\usepackage{amsmath,amssymb,mathtools,unicode-math}
\setmathfont{Latin Modern Math}
\usepackage{enumitem,needspace,xurl,hyperref}
\setlength{\parindent}{0pt}
\setlength{\parskip}{5pt}
\setlength{\emergencystretch}{3em}
\newcommand{\N}{\mathbb{N}}
\newcommand{\Z}{\mathbb{Z}}
\newcommand{\Q}{\mathbb{Q}}
\newcommand{\R}{\mathbb{R}}
\newcommand{\C}{\mathbb{C}}
\newcommand{\F}{\mathbb{F}}
\newcommand{\A}{\mathbb{A}}
\newcommand{\PP}{\mathbb{P}}
\newcommand{\E}{\mathbb{E}}
\newcommand{\eps}{\varepsilon}
\newcommand{\dd}{\,\mathrm{d}}
\begin{document}
\section*{UnsolvedMath ID: 30000918}
\textbf{Source problem code:} catalogue entry
\subsection{Definitions and mathematical statement}
Let $K\subset\C$ be compact and connected with connected complement $\Omega=\C\setminus K$. Let $Y$ be the source's family of Faber-expansion sets $\Gamma\subset\Omega$. For $f\in\mathcal O(\Omega)$ and $\zeta\in\Omega$, its Taylor partial sums are
\[
S_n(f,\zeta)(z)=\sum_{j=0}^{n}\frac{f^{(j)}(\zeta)}{j!}(z-\zeta)^j.
\]
Retain $U(\Omega,\zeta)$ for the source's universal Taylor class and $U(\Omega,\Gamma)$ for its universal Faber class, with their intended allowable approximation sets, target functions and subsequence convergence.

Put
\[
A=\bigcap_{\Gamma\in Y}U(\Omega,\Gamma),\qquad
B=\bigcap_{\zeta\in\Omega}U(\Omega,\zeta).
\]
The source supplies the inclusion $A\subseteq B$ and asks whether $A\subsetneq B$. Equivalently, the exact quantifier pattern is
\[
\exists f\in\mathcal O(\Omega):\quad
\bigl[\forall\zeta\in\Omega,\ f\in U(\Omega,\zeta)\bigr]
\ \text{and}\
\bigl[\exists\Gamma\in Y,\ f\notin U(\Omega,\Gamma)\bigr].
\]
Failure for at least one Faber set suffices; failure for every Faber set is not requested. Intersections of universality classes also do not by themselves require a single subsequence to work simultaneously for every center or expansion set.

The extract does not define admissible Faber-expansion sets, Faber partial sums, approximation compacta/targets or the topology in the two universal classes. These missing definitions prevent replacing the named classes by a unique standard universality convention. The report calls this a doubly connected setting; any nonempty/nondegenerate $K$ assumptions implicit in that description should be confirmed. The original source records a belief in strictness, not a proof or a stronger simultaneous-universality claim.
\subsection{Short English statement}
Is there a holomorphic function universal at every Taylor center that fails universality for at least one admissible Faber-expansion set? The source provides the inclusion and conjectural strictness, but not the detailed Faber family and universality definitions.
\subsection{Sources}
\begin{itemize}
\item Supplied report: \emph{Mini-Workshop: Complex Approximation and Universality} (2008), Oberwolfach Reports, DOI 10.4171/owr/2008/06, pp. 297--346; proposed year 2008. \url{https://doi.org/10.4171/owr/2008/06}.
\item Supplied Vlachou reference: \url{https://doi.org/10.1112/jlms/jdp038}.
\item Local UnsolvedMath registry, permanent ID 30000918, catalogue entry; created 25 August 2026, updated 16 September 2026. Its corrected\_verified note says the domains, classes, inclusion and equivalent existence question are explicit. The original says equality is unknown and strictness is believed in the doubly connected setting. The supplied 21 August 2026 literature triage labels the equality/strictness problem open, distinguishing a simultaneous-Faber theorem for a substantial domain class from the requested comparison. This dated assessment has not been independently verified here.
\end{itemize}
\end{document}
